\chapter{From a calculation to an EBU system}\label{ch:system-design}

\section{The promise and the missing pieces}
Money is useful because it lets people coordinate exchange without re-deriving the full physical history of every purchase. The opening chapters showed why a price alone does not tell us whether the underlying action preserves the conditions of future life. EBU offers a different piece of information: the signed change in a declared measure of physical balance caused by a real transformation. A system built around that information could give people a reason to repair, maintain and regenerate rather than reward only what can be sold now.

The equation is not yet that system. A mathematical result gives a number once state, field and action have been defined. It does not define who may act, whose need has priority, what a person owns, or what happens when a necessary action has negative EBU. Those questions are not embarrassments to be hidden in small print. They are design choices on which the system's usefulness depends.

\section{Five layers that must remain distinguishable}
The first layer is physical description: stocks, flows, carriers, sinks, boundaries and services. Every accepted action has a predecessor and successor state and must pass the relevant conservation and feasibility tests. A useful label is not a substitute for an actual transition.

The second layer is valuation: a declared, versioned potential on the represented state, with its reference values, scales, coupled factors and scientific rationale. Under one fixed field, exact finite EBU is its before/after difference. This layer can be local when affected factors are identifiable; that possibility is important because no society can recalculate a complete planetary state for every small action.

The third layer is attribution. A group of simultaneous actions has one physical successor and one group value. Separate actor records may be useful, but a split of that group value requires its own declared path, rule and meaning. An algebraic decomposition is not automatically a right to spend.

The fourth layer is choice. People and organizations face needs, constraints, histories and incentives. A controller might use an EBU quote to rank physically permissible plans, limit a harmful plan, or make restoration more affordable. It might also protect indispensable service even when the current field gives an adverse sign. The particular policy must be specified before results are observed if it is to be scientifically tested.

The fifth layer is institution: verification, rights, appeal, distribution, public authority and transition from existing money. An accounting unit could coexist with monetary prices, be used only for disclosure, or become part of a more ambitious capacity system. The physical theorem does not select among these arrangements. Their effects on inequality, coercion and access must be examined rather than presumed beneficial.

\section{For what would someone receive EBU?}
This question needs a concrete answer at the level where the term is used. In the physical account, an action receives a signed EBU record when a verified complete transition lowers or raises the declared fixed potential. It is the event, not the actor's virtue, that is measured. If the event lowers the potential, its EBU is positive; if it raises the potential, negative. A zero record can still be a real action.

In a prospective capacity institution, the further question is whether and how that record changes anyone's future ability to authorize actions. One possibility is that verified restoring work creates a bounded claim, while damaging work uses capacity. But a raw identity cannot settle transfers between owners, initial endowments, essential needs, external disturbances or the possibility of cycling through actions as the field changes. A design that ignores these cases could reward strategic damage or strand people whose urgent needs carry negative field value. That is why a ``wallet'' must be a separately justified mechanism with adversarial tests, not a renamed physical receipt.

Imagine a town repairing a leaking water line. The first record is not ``a repair earns EBU.'' It describes the line, the water carrier, any energy and material use, the measured pre-repair loss, and the successor condition of the system. A fixed, declared potential then compares the complete states. If the repair genuinely reduces represented deviation, its physical EBU is positive and can be checked by others. The result gives the town information that a repair invoice alone could not convey. A capacity institution could go further and reward the verified contribution, but it must disclose that extra rule and check that the reward cannot be obtained by first creating the leak. This is the distinction between a useful incentive and a circular credit machine.

The same framework does not need to punish a person for drinking water or receiving care. It can recognize service as a reason to act, physical limits as constraints, and the EBU account as information about consequences. Where service and the current field conflict, an institution has to make its priority explicit. An honest negative account can motivate investment in a better supply route or more resilient infrastructure; hiding the negative number would make those improvements harder to see.

\section{An inspectable event, not an invisible score}
A serious record should say which action occurred; which state and field versions were used; which quantities and carriers changed; how conservation closed; which service it was meant to provide; what evidence supports the measurement; and whether the result was a group value or a specifically defined attribution. Corrections should preserve the original event and explain the new evidence rather than overwrite history.

This information is enough to make a calculation challengeable. It is not a demand for public disclosure of personal life. Privacy, data retention and access are institutional design problems. They should be solved with the minimum measurement needed for the claim, independent review and a way to contest errors.

\section{How an incentive could differ}
Under a price-only decision, maintaining a shared water line may be delayed when the payer cannot capture its broad future benefit. With a trustworthy physical account, the same repair can be described by its effect on represented reliability, water availability and energy use. If an institution then rewards demonstrated restoration or uses that information to finance preventive work, the incentive changes. That is a plausible design pathway, not yet an observed economy-wide outcome.

The design must also resist a tempting shortcut: maximizing the EBU number at every moment. A body stays alive through regulated flows, not by maximizing one hormone or emptying every reserve. A society needs service, reserves, redundancy and adaptation as well as balance. The useful question is whether EBU-guided choices, among real feasible alternatives, help preserve these conditions over time better than relevant comparators. We cannot answer it by looking at one favourable equation.

The next chapter states what kind of proof and evidence would turn that promise into a defensible claim.
